% ═══════════════════════════════════════════════════════════════
% SLIDES - DS-02: Mi nueva casa – Repaso
% Klasse 9, Unidad 2
% ═══════════════════════════════════════════════════════════════

\documentclass[aspectratio=169, 11pt]{beamer}

\usepackage[utf8]{inputenc}
\usepackage[T1]{fontenc}
\usepackage[ngerman]{babel}

\usepackage{helvet}
\usepackage{palatino}
\newcommand{\seriftext}[1]{{\fontfamily{ppl}\selectfont #1}}

\usepackage{tikz}
\usetikzlibrary{calc}
\usepackage{tabularx}
\usepackage{booktabs}
\usepackage{array}

% ═══════════════════════════════════════════════════════════════
% FARBEN
% ═══════════════════════════════════════════════════════════════

\definecolor{spanischrot}{HTML}{C41E3A}
\definecolor{dunkel}{HTML}{222222}
\definecolor{mittelgrau}{HTML}{666666}
\definecolor{hellgrau}{HTML}{F5F5F5}

% ═══════════════════════════════════════════════════════════════
% BEAMER-THEME
% ═══════════════════════════════════════════════════════════════

\usetheme{default}
\usecolortheme{default}

\setbeamercolor{structure}{fg=spanischrot}
\setbeamercolor{frametitle}{fg=dunkel, bg=white}
\setbeamercolor{title}{fg=white}
\setbeamercolor{subtitle}{fg=white}
\setbeamercolor{author}{fg=white}
\setbeamercolor{date}{fg=white}
\setbeamercolor{normal text}{fg=dunkel}

\setbeamerfont{frametitle}{size=\Large, series=\bfseries}
\setbeamerfont{title}{size=\LARGE, series=\bfseries}

\setbeamertemplate{navigation symbols}{}
\setbeamertemplate{footline}{%
    \hfill\textcolor{mittelgrau}{\insertframenumber}\hspace{1em}\vspace{0.5em}
}

% ═══════════════════════════════════════════════════════════════
% BEFEHLE
% ═══════════════════════════════════════════════════════════════

\newcommand{\es}[1]{\seriftext{\textbf{#1}}}
\newcommand{\de}[1]{{\small\textcolor{mittelgrau}{\textit{#1}}}}
\newcommand{\regel}[1]{%
    \tikz[baseline=(X.base)]\node[fill=spanischrot, text=white,
    font=\scriptsize\bfseries, inner sep=3pt, rounded corners=1pt](X){#1};%
}

% ═══════════════════════════════════════════════════════════════
% DOKUMENT
% ═══════════════════════════════════════════════════════════════

\begin{document}

% ═══════════════════════════════════════════════════════════════
% TITELFOLIE
% ═══════════════════════════════════════════════════════════════

{
\setbeamercolor{background canvas}{bg=spanischrot}
\begin{frame}[plain]
    \vfill
    \begin{center}
        {\usebeamerfont{title}\usebeamercolor[fg]{title}Mi nueva casa – Repaso}\\[8pt]
        {\usebeamercolor[fg]{subtitle}Wiederholung + Pretérito perfecto}\\[20pt]
        {\usebeamercolor[fg]{date}\small Spanisch · Klasse 9 · Unidad 2 · DS-02}
    \end{center}
    \vfill
\end{frame}
}

% ═══════════════════════════════════════════════════════════════
% LERNZIELE
% ═══════════════════════════════════════════════════════════════

\begin{frame}{Lernziele}
    \begin{enumerate}
        \item Hausvokabeln und Möbel \textbf{wiederholen}
        \item Bewegungsverben + Präpositionen \textbf{festigen}
        \item Pretérito perfecto mit Umzugsvokabular \textbf{kombinieren}
        \item Einen Umzugsbericht \textbf{schreiben}
    \end{enumerate}

    \vspace{15pt}

    \begin{center}
        \fcolorbox{spanischrot}{hellgrau}{\parbox{0.7\textwidth}{\centering
            \small\textbf{Heute:} Alles zusammenbringen!\\
            Vokabeln + Verben + Perfekt = Umzugsbericht
        }}
    \end{center}
\end{frame}

% ═══════════════════════════════════════════════════════════════
% VOKABEL-QUIZ
% ═══════════════════════════════════════════════════════════════

\begin{frame}{Vokabel-Quiz: ¿Cómo se dice...?}
    \begin{columns}
        \begin{column}{0.5\textwidth}
            \begin{enumerate}
                \item die Küche = ?
                \item das Bett = ?
                \item der Sessel = ?
                \item die Dusche = ?
            \end{enumerate}
        \end{column}
        \begin{column}{0.5\textwidth}
            \begin{enumerate}
                \setcounter{enumi}{4}
                \item der Balkon = ?
                \item der Spiegel = ?
                \item der Schrank = ?
                \item der Flur = ?
            \end{enumerate}
        \end{column}
    \end{columns}

    \vspace{15pt}

    \begin{center}
        \regel{Jetzt} $\rightarrow$ AB Aufgabe 1 (Vokabel-Quiz)
    \end{center}
\end{frame}

% ═══════════════════════════════════════════════════════════════
% VOKABEL-QUIZ LÖSUNG
% ═══════════════════════════════════════════════════════════════

\begin{frame}{Vokabel-Quiz: Solución}
    \begin{columns}
        \begin{column}{0.5\textwidth}
            \begin{enumerate}
                \item \es{la cocina}
                \item \es{la cama}
                \item \es{el sillón}
                \item \es{la ducha}
            \end{enumerate}
        \end{column}
        \begin{column}{0.5\textwidth}
            \begin{enumerate}
                \setcounter{enumi}{4}
                \item \es{el balcón}
                \item \es{el espejo}
                \item \es{el armario}
                \item \es{el pasillo}
            \end{enumerate}
        \end{column}
    \end{columns}

    \vspace{15pt}

    \begin{center}
        {\small Selbstkontrolle: Vergleiche mit AB-01 Kasten V1!}
    \end{center}
\end{frame}

% ═══════════════════════════════════════════════════════════════
% BEWEGUNGSVERBEN
% ═══════════════════════════════════════════════════════════════

\begin{frame}{Los verbos de movimiento + Präpositionen}
    \begin{center}
        \begin{tabular}{cll}
            \textbf{Symbol} & \textbf{Verb} & \textbf{Präposition} \\
            \midrule
            $\downarrow$ & \es{bajar} \de{herunterbringen} & \textbf{de / del} \\[0.5em]
            $\uparrow$ & \es{subir} \de{hochbringen} & \textbf{a / al} \\[0.5em]
            $\oplus$ & \es{poner} \de{stellen} & \textbf{en} \\[0.5em]
            $\ominus$ & \es{sacar} \de{herausnehmen} & \textbf{de / del} \\[0.5em]
            $\rightarrow$ & \es{llevar} \de{hinbringen} & \textbf{a / al} \\[0.5em]
            $\leftarrow$ & \es{traer} \de{herbringen} & \textbf{de / del} \\
        \end{tabular}
    \end{center}

    \vspace{10pt}

    \begin{center}
        \regel{Jetzt} $\rightarrow$ AB Aufgabe 2 (Zuordnung)
    \end{center}
\end{frame}

% ═══════════════════════════════════════════════════════════════
% PRETÉRITO PERFECTO: HABER
% ═══════════════════════════════════════════════════════════════

\begin{frame}{El pretérito perfecto: Konjugation von \es{haber}}
    \begin{columns}
        \begin{column}{0.5\textwidth}
            \begin{center}
                \begin{tabular}{ll}
                    yo & \es{he} \\
                    tú & \es{has} \\
                    él/ella/usted & \es{ha} \\
                    nosotros/-as & \es{hemos} \\
                    vosotros/-as & \es{habéis} \\
                    ellos/ellas/ustedes & \es{han} \\
                \end{tabular}
            \end{center}
        \end{column}
        \begin{column}{0.5\textwidth}
            \begin{center}
                \fcolorbox{spanischrot}{hellgrau}{\parbox{0.85\textwidth}{\centering
                    \textbf{Bildung:}\\[4pt]
                    \es{haber} + Partizip\\[4pt]
                    \es{He puesto la lámpara.}
                }}
            \end{center}
        \end{column}
    \end{columns}

    \vspace{15pt}

    \de{Du kennst das schon aus Unidad 1!}
\end{frame}

% ═══════════════════════════════════════════════════════════════
% PRETÉRITO PERFECTO: PARTIZIPIEN
% ═══════════════════════════════════════════════════════════════

\begin{frame}{El pretérito perfecto: Partizipien}
    \begin{columns}
        \begin{column}{0.5\textwidth}
            \textbf{Regelmäßig:}

            \vspace{8pt}

            \begin{tabular}{lll}
                \textbf{-ar} & hablar & habl\textcolor{spanischrot}{\textbf{ado}} \\[0.4em]
                \textbf{-er} & comer & com\textcolor{spanischrot}{\textbf{ido}} \\[0.4em]
                \textbf{-ir} & vivir & viv\textcolor{spanischrot}{\textbf{ido}} \\
            \end{tabular}
        \end{column}
        \begin{column}{0.5\textwidth}
            \textbf{Unregelmäßig:}

            \vspace{8pt}

            \begin{tabular}{ll}
                poner & \textcolor{spanischrot}{\textbf{puesto}} \\[0.4em]
                hacer & \textcolor{spanischrot}{\textbf{hecho}} \\
            \end{tabular}
        \end{column}
    \end{columns}

    \vspace{15pt}

    \begin{center}
        \regel{Jetzt} $\rightarrow$ AB Kasten G1 ergänzen
    \end{center}
\end{frame}

% ═══════════════════════════════════════════════════════════════
% REDEMITTEL
% ═══════════════════════════════════════════════════════════════

\begin{frame}{Redemittel: Vergangene Handlungen}
    \begin{center}
        \begin{tabular}{ll}
            \es{He puesto el/la... en...} & \de{Ich habe ... in ... gestellt.} \\[0.6em]
            \es{Hemos subido los/las... al piso.} & \de{Wir haben ... hochgebracht.} \\[0.6em]
            \es{Mi padre ha sacado...} & \de{Mein Vater hat ... herausgenommen.} \\[0.6em]
            \es{¿Qué has hecho?} & \de{Was hast du gemacht?} \\
        \end{tabular}
    \end{center}

    \vspace{10pt}

    \textbf{Zeitangaben:}\\
    \es{Esta semana... / Hoy... / Esta mañana...}

    \vspace{10pt}

    \begin{center}
        \regel{Jetzt} $\rightarrow$ AB Kasten R1 lesen
    \end{center}
\end{frame}

% ═══════════════════════════════════════════════════════════════
% KOMBINATION
% ═══════════════════════════════════════════════════════════════

\begin{frame}{Kombination: Perfekt + Objektpronomen}
    \begin{center}
        \fcolorbox{spanischrot}{hellgrau}{\parbox{0.8\textwidth}{\centering
            \vspace{8pt}
            \es{He puesto la lámpara en el salón.}\\
            $\downarrow$\\
            \es{\textcolor{spanischrot}{La} he puesto en el salón.}\\
            \vspace{8pt}
        }}
    \end{center}

    \vspace{15pt}

    \begin{center}
        \regel{Wichtig} Das Pronomen steht \textbf{VOR} \es{haber}!\\[8pt]
        \es{\textcolor{spanischrot}{Lo} he llevado...} \quad
        \es{\textcolor{spanischrot}{La} he puesto...} \quad
        \es{\textcolor{spanischrot}{Los} hemos subido...}
    \end{center}
\end{frame}

% ═══════════════════════════════════════════════════════════════
% LÜCKENTEXT
% ═══════════════════════════════════════════════════════════════

\begin{frame}{Übung: Lückentext}
    \begin{center}
        {\small Ergänze mit dem pretérito perfecto und Objektpronomen:}
    \end{center}

    \vspace{10pt}

    \begin{enumerate}
        \item Yo \_\_\_\_\_\_\_\_ (poner) mi cama en el dormitorio.
        \item Mi padre \_\_\_\_\_\_\_\_ (subir) las cajas al piso.
        \item La lámpara – \_\_ \_\_\_\_\_\_\_\_ (poner, yo) al lado de la cama.
        \item ¿Y el sofá? Mi hermano \_\_ \_\_\_\_\_\_\_\_ (llevar) al salón.
    \end{enumerate}

    \vspace{10pt}

    \begin{center}
        \regel{Jetzt} $\rightarrow$ AB Aufgabe 3 (Lückentext)
    \end{center}
\end{frame}

% ═══════════════════════════════════════════════════════════════
% SÄTZE BILDEN
% ═══════════════════════════════════════════════════════════════

\begin{frame}{Übung: Sätze bilden}
    \begin{center}
        \textbf{Modell:}\\[6pt]
        Clara / poner / la lámpara / en el salón\\
        $\downarrow$\\
        \es{Clara ha puesto la lámpara en el salón.}
    \end{center}

    \vspace{15pt}

    \begin{enumerate}
        \item Nosotros / subir / los muebles / al piso
        \item Papá / sacar / la nevera / del camión
        \item Yo / llevar / el televisor / al dormitorio
        \item Mamá / poner / el cuadro / en la pared
    \end{enumerate}

    \vspace{10pt}

    \begin{center}
        \regel{Jetzt} $\rightarrow$ AB Aufgabe 4
    \end{center}
\end{frame}

% ═══════════════════════════════════════════════════════════════
% TEXTPRODUKTION
% ═══════════════════════════════════════════════════════════════

\begin{frame}{Textproduktion: Tu informe de la mudanza}
    \begin{center}
        Schreibe deinen eigenen \textbf{Umzugsbericht}!
    \end{center}

    \vspace{10pt}

    \textbf{Anforderungen:}
    \begin{itemize}
        \item Mindestens \textbf{6 Sätze} im pretérito perfecto
        \item Verwende verschiedene \textbf{Bewegungsverben}
        \item Nutze \textbf{Zeitangaben} (Esta semana, Hoy...)
    \end{itemize}

    \vspace{10pt}

    \textbf{Hilfen:}\\
    \es{Esta semana mi familia y yo nos hemos mudado...}\\
    \es{He puesto... / Mi padre ha subido... / Hemos trabajado mucho.}

    \vspace{10pt}

    \begin{center}
        \regel{Jetzt} $\rightarrow$ AB Aufgabe 5 (min. 6 Sätze)
    \end{center}
\end{frame}

% ═══════════════════════════════════════════════════════════════
% ZUSAMMENFASSUNG
% ═══════════════════════════════════════════════════════════════

\begin{frame}{Zusammenfassung}
    \begin{columns}
        \begin{column}{0.33\textwidth}
            \textbf{Wiederholung:}
            \begin{itemize}
                \item Hausvokabeln
                \item 6 Bewegungsverben
                \item Präpositionen
                \item Objektpronomen
            \end{itemize}
        \end{column}
        \begin{column}{0.33\textwidth}
            \textbf{Neu kombiniert:}
            \begin{itemize}
                \item Pretérito perfecto
                \item Pronomen vor \es{haber}
                \item \es{Lo he puesto...}
            \end{itemize}
        \end{column}
        \begin{column}{0.33\textwidth}
            \textbf{Kommunikation:}
            \begin{itemize}
                \item Umzugsbericht
                \item Vergangenes beschreiben
            \end{itemize}
        \end{column}
    \end{columns}

    \vspace{15pt}

    \begin{center}
        \fcolorbox{spanischrot}{hellgrau}{\parbox{0.6\textwidth}{\centering
            \textbf{Nächste Stunde:}\\
            Klausurvorbereitung U1 + U2
        }}
    \end{center}
\end{frame}

% ═══════════════════════════════════════════════════════════════
% ABSCHLUSS
% ═══════════════════════════════════════════════════════════════

{
\setbeamercolor{background canvas}{bg=spanischrot}
\begin{frame}[plain]
    \vfill
    \begin{center}
        {\usebeamerfont{title}\usebeamercolor[fg]{title}¡Buen trabajo!}\\[20pt]
        {\usebeamercolor[fg]{subtitle}\small Hausaufgabe: AB Aufgabe 5 fertigstellen}
    \end{center}
    \vfill
\end{frame}
}

\end{document}
